\section{Panorama general}
\label{sec:panorama}

\subsection{De qu\'e se trata el proyecto}

El objetivo es construir un sismógrafo multicanal de bajo costo para
caracterizaci\'on geot\'ecnica del suelo mediante el m\'etodo MASW
(\emph{Multichannel Analysis of Surface Waves}): se golpea el suelo con una
fuente de impacto, se registra la llegada de las ondas Rayleigh en un arreglo
lineal de ge\'ofonos, se construye la imagen de dispersi\'on
velocidad-de-fase/frecuencia y se invierte esa curva para obtener el perfil de
velocidad de onda de corte $V_S(z)$.

Lo que distingue a este sistema de un sism\'ografo comercial es el conjunto de
restricciones autoimpuestas:

\begin{itemize}[itemsep=0.2em]
\item \textbf{El transductor es el que hay}, no el que conviene: ge\'ofonos
SM-24 de \SI{10}{\hertz} disponibles en la Facultad. Para llegar a las
profundidades de inter\'es hace falta informaci\'on por debajo de su frecuencia
natural, de modo que \emph{extender la banda hacia abajo por v\'ia electr\'onica
deja de ser un refinamiento y pasa a ser una condici\'on de dise\~no}. La
extensi\'on efectiva todav\'ia debe medirse de extremo a extremo.
\item \textbf{Nodos inal\'ambricos e independientes}, no un cable de decenas de
metros con un digitalizador central. Eso convierte la sincronizaci\'on entre
canales en un problema de dise\~no de primer orden.
\item \textbf{Orquestaci\'on y revisi\'on en herramientas propias}, desde la
adquisici\'on hasta una fachada de inversi\'on con motores propios y de terceros.
La capacidad multimodal existe, pero el estado persistido usa s\'olo una curva
etiquetada como modo~0, sin identificaci\'on independiente como fundamental.
\end{itemize}

\subsection{Resumen autoritativo al 12 de agosto de 2026}

Esta tabla es el punto de entrada para resumir el proyecto. Ante una diferencia
con un relato cronol\'ogico posterior, prevalecen los archivos de resultados, el
firmware vigente y los metadatos indicados aqu\'i.

\begin{center}\footnotesize
\begin{tabularx}{\textwidth}{@{}l X@{}}
\toprule
\textbf{Eje} & \textbf{Estado que puede afirmarse} \\
\midrule
Transductor & SM-24 disponible en la Facultad: $f_n=\SI{10}{\hertz}$,
sensibilidad nominal \SI{28.8}{\volt\second\per\meter}, $\zeta_0=0{,}25$,
bobina \SI{375}{\ohm} y masa m\'ovil \SI{11}{\gram}; la hoja de datos no
especifica inductancia. \\
\addlinespace[2pt]
Compensaci\'on & En la tanda calibrada, $\zeta_0=0{,}25$ es el objetivo del
numerador y $\zeta_1=937{,}396$ el denominador BP. $83{,}661$ pertenece a la
tanda descartada. El polo inferior ideal es \SI{5.44}{\milli\hertz}, pero la
muesca medida qued\'o \SI{26.33}{\decibel} menos profunda que la de dise\~no. \\
\addlinespace[2pt]
Modelo del operacional & Los scripts combinaron $A_0=\SI{90}{\decibel}$ con
\SI{8}{\mega\hertz} como polo dominante. Si este valor es GBW, el polo monopolo
es aproximadamente \SI{253}{\hertz}. La relaci\'on ya se corrigi\'o en los scripts;
la s\'intesis y el contraste deben repetirse y el ajuste previo no se usa como validaci\'on. \\
\addlinespace[2pt]
DA-AFE & Sobre dos placas y ocho configuraciones, la salida LP pas\'o de
\SIrange{2.9}{45.8}{\percent} a \SIrange{0.4}{3.9}{\percent} del fondo de
escala, con 10 reinicios apareados por configuraci\'on y un \'unico PI. \\
\addlinespace[2pt]
ADC y adquisici\'on & Firmware actual: \SI{2604}{\hertz} nativos, 18 bits,
$F_s=2604/N$ con $N=1,\ldots,100$ y cuatro rangos. Las tasas 2929 y
\SI{1020}{\hertz} son hist\'oricas de campo. La ruta cr\'itica usa DFB, DMA y
\texttt{superMaquina} en UDB. \\
\addlinespace[2pt]
Sincronizaci\'on & PRESTART/HOT\_WAIT asegura disponibilidad; el
\texttt{CMD\_START} llega por ESP-NOW y cada esclavo genera el flanco local al
PSoC. La observaci\'on $<\SI{400}{\micro\second}$ a unos \SI{5}{\meter} es
preliminar y anterior a la medici\'on completa de la topolog\'ia vigente. \\
\addlinespace[2pt]
Campo & Apertura sint\'etica: martillo fijo en $x=0$ y un ge\'ofono m\'ovil. El
grupo principal tiene 21 posiciones entre \SI{10}{} y \SI{50}{\meter},
$\Delta x=\SI{2}{\meter}$ y $L=\SI{40}{\meter}$. El orden de profundidad con
soporte conservador es \SI{20}{\meter}, no \SI{50}{\meter}. \\
\addlinespace[2pt]
MASW & El exportado de julio tiene 51 puntos y desajuste declarado de
\SI{2.554}{\percent}, pero llega a $\lambda=\SI{130.6}{\meter}$ y queda fuera de
la evidencia. El estado posterior persistido usa $\lambda=\SIrange{4.001}{26.732}{\meter}$,
\texttt{maswavespy(port)}, modo~0, siete capas y desajuste \SI{7.1063}{\percent}. No hay
validaci\'on SPT ni resultado multimodal comparado. \\
\addlinespace[2pt]
Pendientes de cierre & Modelo no ideal del operacional, banda compensada medida, \emph{jitter} extremo a
extremo, apareamiento entre canales, mayor apertura y energ\'ia de fuente,
ruido referido a entrada y contraste geot\'ecnico externo. \\
\bottomrule
\end{tabularx}
\end{center}

\begin{figure}[!ht]
\centering
\includegraphics[width=\textwidth]{arquitectura_vigente.png}
\caption{Arquitectura vigente y frontera entre el camino determinista de
muestra y las funciones de calibraci\'on, comunicaci\'on, almacenamiento y
procesamiento.}
\label{fig:arquitectura-vigente-hist}
\end{figure}

\subsection{Escala del trabajo}

\begin{figure}[!ht]
\centering
\includegraphics[width=\textwidth]{timeline_commits.png}
\caption{Actividad de desarrollo semanal en los diez repositorios del proyecto.
La intensidad codifica cantidad de \emph{commits} por semana en escala
comprimida (ra\'iz cuadrada), porque de otro modo la sesi\'on de 233
\emph{commits} de abril aplastar\'ia todo el resto. La fila
\emph{Superproyecto} concentra la historia anterior a la modularizaci\'on de
julio de 2026: hasta ese momento todo viv\'ia en un \'unico repositorio, y por
eso las filas de los subm\'odulos arrancan m\'as tarde aunque el trabajo sea
anterior.}
\label{fig:timeline}
\end{figure}

\begin{center}\small
\begin{tabularx}{\textwidth}{@{}lX r@{}}
\toprule
\textbf{Dimensi\'on} & \textbf{Detalle} & \textbf{Cantidad} \\
\midrule
Per\'iodo activo & 25 feb 2026 --- 12 ago 2026 & 25 semanas \\
\emph{Commits} & suma reconstruida de los diez repositorios & \textasciitilde{}1261 \\
D\'ias con actividad & registrados en la bit\'acora & 66 \\
Bit\'acora & entradas diarias con an\'alisis de \emph{diff} & \textasciitilde{}1{,}3~MB \\
Bibliograf\'ia & base curada con DOI y relevancia & 65 papers \\
Notas de concepto & vault Obsidian, conceptos at\'omicos & \textasciitilde{}153 \\
Datos crudos de campo & tres sitios, LFS con dedup SHA-256 & \textasciitilde{}29~GiB \\
Datos procesados & \emph{waterfalls}, dispersi\'on, perfiles & \textasciitilde{}1{,}1~GiB \\
\bottomrule
\end{tabularx}
\end{center}

\subsection{Los diez repositorios}

En julio de 2026 el monolito se parti\'o en un superproyecto con subm\'odulos
organizados \emph{por prop\'osito}, no por lenguaje. La organizaci\'on final es:

\begin{center}\small
\begin{tabularx}{\textwidth}{@{}l l X@{}}
\toprule
\textbf{Ruta} & \textbf{Repositorio} & \textbf{Contenido} \\
\midrule
\texttt{src/firmware/psoc} & Tesis-firmware-psoc & Adquisici\'on: cadena anal\'ogica configurable, ADC $\Delta\Sigma$, DFB, DMA, m\'aquina de estados en UDB, FatFs sobre SD \\
\texttt{src/firmware/esp32} & Tesis-firmware-esp32 & Red ESP-NOW maestro/esclavo, sincronizaci\'on de disparo, volcado de capturas, servidor web embebido \\
\texttt{src/interfaces/python} & Tesis-interfaces-python & Servidor FastAPI, interfaz web de revisi\'on de campo, cola de an\'alisis \\
\texttt{src/interfaces/matlab} & Tesis-interfaces-matlab & Osciloscopio virtual y bancos de prueba contra el hardware \\
\texttt{src/calculos\_modelados/matlab} & Tesis-calculos-matlab & Modelo del SM-24, compensadores, identificaci\'on de la cadena anal\'ogica, martillo de leva en Simscape \\
\texttt{src/calculos\_modelados/python} & Tesis-calculos-python & C\'alculo de compensadores, MASW, martillo de leva 2D \\
\texttt{docs} & Tesis-documentacion & Vault de investigaci\'on, bit\'acora, papers, art\'iculo URUCON, esta presentaci\'on \\
\texttt{data} & Tesis-datos & Registros crudos y procesados de campo y de osciloscopio \\
\texttt{PCBs} & Tesis-PCBs & Dise\~nos de placa \\
\midrule
\multicolumn{3}{@{}l@{}}{\emph{(el superproyecto \texttt{Tesis} sostiene los \emph{gitlinks} y la historia previa a la partici\'on)}} \\
\bottomrule
\end{tabularx}
\end{center}

\subsection{Las etapas}

\begin{center}\small
\begin{tabularx}{\textwidth}{@{}l l X@{}}
\toprule
\textbf{\S} & \textbf{Per\'iodo} & \textbf{Qu\'e pas\'o} \\
\midrule
\ref{sec:et1} & 25 feb -- 23 mar & Marco te\'orico: Foti completo, 65 papers, scripts de validaci\'on \\
\ref{sec:et2} & 10 -- 15 abr & Investigaci\'on asistida: 140 conceptos at\'omicos \\
\ref{sec:et3} & 17 abr -- 20 may & Electr\'onica anal\'ogica: modelo del SM-24, compensaci\'on, dos pivotes de topolog\'ia \\
\ref{sec:et4} & 24 may -- 3 jun & Red ESP-NOW: sincronizaci\'on de disparo entre nodos \\
\ref{sec:et5} & 4 -- 10 jun & Consolidaci\'on del nodo: firmware \'unico e interfaz web embebida \\
\ref{sec:et6} & 14 jun -- 2 jul & Auto-calibraci\'on, DMA y \texttt{superMaquina} \\
\ref{sec:et7} & 3 -- 9 jul & Campa\~nas de campo y \emph{pipeline} MASW en Python \\
\ref{sec:et8} & 11 -- 12 jul & Auditor\'ia pre-campo, SD en el PSoC, tres bugs cr\'iticos \\
\ref{sec:et9} & 15 -- 23 jul & Art\'iculo URUCON y caracterizaci\'on medida de la cadena \\
\ref{sec:et10} & 21 jul -- 1 ago & Ruptura arquitect\'onica y portaci\'on de la interfaz a web \\
\ref{sec:et11} & 2 -- 12 ago & Placa nueva: rama \texttt{cambios-hardware} \\
\bottomrule
\end{tabularx}
\end{center}
